\section{Conclusion and Future Work}
\label{sec:conclusion}

This work investigated whether a small set of modular skills can transform the solutions generated by an evolving population into a reusable memory that mitigates catastrophic forgetting. To address this question, the Evolutionary Neuroplastic Agent (ENA) was introduced. Across the evaluated CartPole sequences, the empirical ablation supports this architecture. Only variants with failure-driven search master the more demanding environments, requiring increased environment interaction. Among these variants, only the architecture that indexes its archive by detected physical regimes preserves acquired competencies across switches. Introducing the archive also alters test-time policy selection, an effect that the current ablation does not decouple. Policy retrieval via Bayesian policy reuse proved significantly more reliable than relying solely on change detection. The memory is constructed online without external supervision, using detection mechanisms tailored to each domain. Whether dynamically modulating search effort based on failure outperforms uniform effort allocation remains an open question. Single-policy learners trained on identical sequences suffer catastrophic forgetting of previously acquired behaviors. Conversely, they highlight regimes where ENA is less competitive: on CartPole, adaptation is slower than PPO, and on HalfCheetah, asymptotic returns are lower than those of SAC. Both limitations align with known properties of evolutionary search, which exhibits high sample complexity and scales unfavorably with parameter count~\cite{khadka2018evolutionary,vemula2019contrasting}.

The most direct next step is a hybrid of ENA with gradient-based learners. Evolution and gradient descent can be combined in multiple configurations, offering several promising avenues. Gradient learners such as SAC or PPO could serve as local optimizers within each niche to accelerate adaptation and achieve higher asymptotic returns. ENA's environment detection, Plasticity, and physics-indexed Hall of Fame would preserve niche-specific competencies across dynamic switches. Alternatively, the architectures could partition the policy network—with gradient descent optimizing high-dimensional feature representations and evolutionary search tuning low-dimensional control layers—or alternate sequentially. The evolutionary population could also supply diverse rollouts to off-policy replay buffers, as in evolutionary reinforcement learning~\cite{khadka2018evolutionary}. A hybrid architecture would furthermore provide a direct comparison against continual-RL frameworks designed for contextual non-stationarity: OWL~\cite{kessler2022same}, RL-CD~\cite{dasilva2006dealing}, MBCD~\cite{alegre2021minimum}, SWOKS~\cite{dick2024statistical}, and TeLAPA~\cite{lillo2026beyond}.

The modular skills may also transfer directly to those frameworks. OWL and TeLAPA depend on explicit task boundaries; their authors identify automated task boundary detection during training~\cite{kessler2022same} and handling ambiguous transitions~\cite{lillo2026beyond} as primary challenges. ITE initiates adaptation based on a predefined performance threshold and maintains no memory across sequential conditions~\cite{cully2015robots}. ENA's environment detection identifies transitions and recognizes recurring regimes directly from unassisted observations, providing capabilities that existing methods either assume or lack. Furthermore, a physics-indexed archive could preserve adapted policies for deployment when previously encountered conditions recur.

Several further directions remain open. Repeated recurrence during training and randomized task orders would test the memory beyond the evidence here. For systems that cannot be reset, the probe could be superimposed on the agent's own actions as a small dither. The contributions of the belief's two sources of evidence, the dynamics likelihoods and the test returns, remain to be separated. Specialist selection could also be learned, as a contextual bandit over the detector's evidence. Finally, the population's evaluations are independent and parallelize naturally on accelerators, which would allow larger populations and longer non-stationary streams.
